\documentclass[journal=jcisd8,manuscript=article]{achemso}
\usepackage{fontspec}
\setmainfont{DejaVu Serif}
\setsansfont{DejaVu Sans}
\setmonofont{DejaVu Sans Mono}[Scale=0.85]
\usepackage{graphicx}
\usepackage{tabularx}
\usepackage{caption}
\captionsetup{singlelinecheck=false}
\usepackage{url}
\usepackage[hyphens]{xurl}
\sloppy
\author{[TODO: authors]}
\affiliation{[TODO: affiliation]}
\email{[TODO: corresponding author email]}
\title{A validated, reproducible pipeline for screening graftable redox motifs for non-aqueous organic redox flow batteries}
\begin{document}
\begin{tocentry}
\includegraphics[width=3.25in,height=1.75in,keepaspectratio]{figures/toc_graphic.png}
\end{tocentry}
\begin{abstract}
Choosing a redox-active molecule to attach (graft) onto a polymer for a non-aqueous organic redox flow battery needs three things. Its properties must be computed for the molecule as it is actually attached, each calculation must be checked automatically, and the method must be tested against measured values. Computational screens rarely provide all three. We present a reproducible pipeline that models each molecule with a small cap standing in for the polymer. It computes redox potentials (vs ferrocene, in acetonitrile), capacity, and reorganization energies with density functional theory (DFT), and checks every redox state before ranking. Tested without any fitting on 39 measured potentials (30 molecules, 7 studies), the computed potentials have a mean absolute error (MAE) of 0.247 V, most of it a shared offset of +0.19 V, so rankings are more reliable than absolute values. Applied to 91 candidates, the pipeline shows that the options are dominated by anolytes (85 anolytes, 6 catholytes). Grafted 5,10-dihydrophenazines fill part of the empty catholyte region. A series that changes one substituent at a time gives a simple rule: electron-withdrawing groups on the nitrogen not attached to the polymer raise the potential, and electron-donating groups on the ring lower it. Pairing these catholytes with anolytes from the pool gives proposed cells of about 1.7–1.9 V after correcting for each family's measured error. The code, worked examples, and tests are released.
\end{abstract}
\section{Introduction}

Non-aqueous organic redox flow batteries (ORFBs) are an attractive route to grid-scale electrochemical storage. They decouple energy from power, reach cell voltages well beyond the stability window of water, and are built from abundant, tunable organic redox couples rather than scarce transition metals \cite{soloveichik2015}. One of the more promising device architectures immobilizes the redox-active species on a soluble or swellable polymer, so that the active material is retained by a size-selective membrane and crossover is suppressed \cite{janoschka2015}. In practice these materials are most easily made not by polymerizing a custom redox monomer but by grafting a ready-made redox molecule onto a reactive polymer. The usual choice is chloromethylated (Merrifield-type) polystyrene, whose benzyl chloride groups are alkylated by a nucleophilic site on the redox molecule \cite{merrifield1963}. This grafting route decouples the chemistry of the backbone from the chemistry of the redox center and makes a large library of candidate motifs synthetically plausible from a single polymer platform. The design space it opens, however, is badly unbalanced. The catalog of well-behaved organic anolytes — n-type couples based on quinones, imides, and related electron-poor scaffolds — is deep and still growing \cite{ding2018}, whereas robust, high-potential catholytes remain scarce \cite{kwon2019}. Because a flow cell needs both a catholyte and an anolyte, the shortage of good catholytes limits which cells can be built. This is the kind of gap that a well-designed computational screen should be able to map and help fill.

Computation is in principle well suited to this task. Redox potentials, and the reorganization energies that govern electron-transfer rates, can be computed from electronic-structure calculations in solvent, and density functional theory (DFT) with an implicit solvent model has been used widely to rank organic redox couples \cite{marenich2014,er2015}. In practice, though, most published screens are difficult to trust or to reuse for the specific problem of grafted motifs. Three shortcomings recur. First, the molecular model is often chosen for convenience rather than to match the real material. A bare parent heterocycle, an N–H or N-methyl analogue, or a soluble monomer stands in for the grafted species, even though the attachment changes the electronic environment of the redox center and therefore its potential. A capped graft, a soluble parent, and an oligomer are genuinely different models and should not be conflated. Second, reliability is rarely gated per candidate. Spin contamination, unconverged or distorted geometries, missing charge states, and inconsistent reference points can silently corrupt a reported potential, and screens seldom document a per-state quality-control (QC) policy or state which candidates were held back because they failed it. Third, and most importantly, computed numbers are frequently reported without an external calibration against measured potentials under matched conditions, so the reader cannot tell whether a ranking reflects chemistry or method artifacts. Underlying all three is a reproducibility problem. The level of theory, solvent, standard state, reference electrode, and software versions are often not fully stated, and the raw records behind a figure are not released, so neither the ranking nor its error bars can be rebuilt independently \cite{lejaeghere2016}. A computed property, or a completed calculation, is not by itself evidence of synthetic feasibility or of material performance; a screen is only as useful as the discipline with which these distinctions are maintained.

This work addresses those gaps with a single pipeline for screening graftable redox molecules for non-aqueous ORFBs. Each design choice answers one of the problems above. (i) Each molecule is modeled as it is attached: the redox center carries a 4-methylbenzyl group that stands in for the polystyrene backbone, and each charge state is built and optimized on its own (Figure 1). (ii) Every redox state is checked automatically (structure, spin, protocol, completeness), and states that fail are kept with their reason rather than silently dropped. (iii) All potentials are computed at one level of theory (ωB97M-V/def2-TZVPD in implicit acetonitrile) and reported against ferrocene (Fc). They are tested, without any fitting, against measured potentials: the mean absolute error is 0.247 V, most of it a shared offset of +0.19 V, so rankings are more reliable than absolute values. (iv) Candidates are ranked on Pareto fronts of potential and capacity, that is, the set of candidates that no other candidate beats on both properties at once; anolytes and catholytes are ranked separately, with the potential uncertainty shown and no weighted-sum score. Each property comes with its own level of confidence: the outer-sphere reorganization energy is a continuum estimate, while the inner-sphere term is treated as qualitative. (v) The code is released with worked examples that reproduce stored results byte for byte, a large test suite, and records that tie each reported number to the calculations it came from.

\begin{figure}[htbp]
\centering
\includegraphics[width=0.98\textwidth]{figures/fig1_grafting.png}
\caption{\textbf{A redox motif is modeled as an exact capped graft, with redox states defined per motif. (a) Grafting abstraction.} Each candidate is built by attaching a redox motif (here 5,10-dihydrophenazine, DMPZ, the p-type catholyte scaffold used in the showcase; both ring nitrogens are N–H in the free motif) to a \emph{4-methylbenzyl} cap that serves as a Merrifield (poly(vinylbenzyl)) backbone proxy; the ring methyl \emph{para} to the benzylic linker stands in for the polymer chain. Grafting is a balanced local substitution onto 4-methylbenzyl chloride (\texttt{Cc1ccc(CCl)cc1}), releasing HCl. \textbf{Only one nitrogen (N5) carries the Merrifield backbone graft} (highlighted orange: the new benzylic C–N bond); \textbf{the other nitrogen (N10) carries a small ethyl \emph{tuning} cap, not a second graft.} The depicted capped graft, \emph{N}-(4-methylbenzyl)-\emph{N}′-ethyl-DMPZ (\texttt{CCN1c2ccccc2N(Cc2ccc(C)cc2)c2ccccc21}), is identical to the registered model (\texttt{config/custom\_candidates.py}, id \texttt{exp\_dmpz\_mb\_ethyl}). \textbf{(b) Per-motif redox states.} The redox states are defined on the capped graft as adjacent one-electron couples from the resting state; for this p-type motif, neutral → ox1 (+1, doublet radical cation) → ox2 (+2, singlet dication), reported vs Fc/Fc⁺. Because \textbf{both nitrogens are substituted (N,N′-disubstituted), N–H deprotonation on oxidation is blocked, giving a reversible 2e⁻ couple} \cite{kwon2019} — the design basis for using DMPZ as a 2-electron catholyte. (n-type anolyte motifs instead reduce; ambipolar motifs carry couples on both sides.)}
\label{fig:1}
\end{figure}

We demonstrate the pipeline on a pool of 91 graftable candidates (the \emph{pool}) from the quinone, imide, viologen, dihydrophenazine, and phenothiazine families. The screen confirms quantitatively that the accessible space is strongly anolyte-dominated (85 anolyte versus 6 catholyte candidates), and the sparsity of the catholyte Pareto front motivates a focused, matched-substituent study of grafted 5,10-dihydrophenazine two-electron catholytes. We call this grafted scaffold DMPZ; in the literature DMPZ usually denotes the 5,10-dimethyl compound, whereas here one nitrogen carries the polymer cap and the other (N′) carries a small substituent. From series that differ by a single substituent we extract a simple design rule: electron-withdrawing N′-substituents raise the potential, and electron-donating ring groups lower it. This rule places candidates in the previously empty catholyte region, at computed potentials of +0.2 to +0.6 V versus Fc (probably about 0.3 V lower in experiment; the +0.2 V member was added after the 91-candidate pool was fixed). Pairing these catholytes with anolytes from the pool suggests non-aqueous full cells of about 1.7–1.9 V; we stress that these are computational proposals, not demonstrations of capacity, cycling, crossover, or synthetic feasibility. The remainder of the paper is organized as follows. We first describe the grafted-molecule model and the nine-stage pipeline, including its two levels of thermal correction (a cheap one for the whole pool and a more expensive one for confirming leads) and its automatic quality checks. We then present the test against experiment, apply the pipeline to the 91-candidate pool and its Pareto fronts, and finally present the DMPZ catholyte showcase and the proposed full cells, before closing with the method's limitations and outlook.

\section{Computational Methods}

We screen graftable redox molecules with one reproducible driver (\texttt{redox.workflows.motifs.commands}). It takes each listed molecule through nine stages (Figure 2): (1) build the structure; (2) pre-optimize it with a machine-learned interatomic potential (MLIP); (3) optimize each redox state with DFT in implicit solvent; (4) compute refined single-point energies; (5) add thermal corrections; (6) derive the redox properties; (7) run the quality checks; (8) rank the candidates; and (9) audit the exported tables. Every result is stored under a hash of its inputs, so a change in geometry, method, charge, or spin creates a new record instead of overwriting an old one. Table 1 collects the level of theory and the conventions used throughout. The Supporting Information (SI) gives the equations used for each property and a map of which code module carries out each stage.

\begin{figure}[htbp]
\centering
\includegraphics[width=0.98\textwidth]{figures/fig2_pipeline.png}
\caption{\textbf{A single reproducible driver takes each grafted motif from structure to quality-checked redox properties in nine stages.} The driver \texttt{redox.workflows.motifs.commands()} runs: (1) build the grafted-motif model (RDKit ETKDG ensemble, 4-methylbenzyl cap, redox states per motif); (2) UMA MLIP gas-phase pre-optimization (charge-/spin-aware seed); (3) DFT geometry optimization in SMD(MeCN) for each state [neutral / ox1 / ox2] at r²SCAN-D4/def2-SVP(D) with a GFN2-xTB ALPB(MeCN) warm start, the three states run concurrently; (4) refined single-point energies at ωB97M-V(VV10)/def2-TZVPD in SMD(MeCN) using one diffuse basis for every charge state and for the Fc reference; (5) two-tier thermochemistry; (6) derived properties kept as separate quantities; (7) QC coverage and cross-point consistency; (8) Pareto ranking; and (9) an independent arithmetic audit of the exported tables. Each stage is addressed by content hash, so a change in geometry, protocol, or (charge, multiplicity) resolves a new record rather than overwriting an old one. Boxes are coloured by compute resource (legend): GPU/DFT, GPU/ML (UMA), CPU, and QC/scoring. The inset tabulates the level of theory and standard-state conventions (manuscript Table 1). External experimental validation is not shown here; it is Figure 4.}
\label{fig:2}
\end{figure}

\subsection{Grafted-motif model}

Each candidate is a redox-active molecule (a \emph{motif}) attached to a single \emph{4-methylbenzyl} cap, which stands in for a Merrifield-type polystyrene backbone: the ring methyl \emph{para} to the attachment point represents the polymer chain. The graft is built by substituting the motif onto 4-methylbenzyl chloride, keeping atoms and charge balanced. Two reaction types are used: N-alkylation of a pyridine-type nitrogen (Cl⁻ leaves) and substitution at an N–H or O–H group (HCl leaves, which in practice needs a base). No base, counterion, or reaction rate is modeled. Starting from the molecule as prepared, the redox states are the successive one-electron steps. \emph{p-type} motifs, used as catholytes, are oxidized (neutral → +1 → +2); \emph{n-type} motifs, used as anolytes, are reduced (neutral → −1 → −2, or +2 → +1 → 0 for viologens); ambipolar motifs do both. A \textbf{capped graft, a soluble parent, and an oligomer are three distinct models}: the capped graft fixes the local electronic structure and the attachment chemistry, but not chain packing, polymer solubility, or interactions between neighboring grafts. For capacity, the cap's methyl is replaced by one backbone carbon per graft site (CH₃–Ar → –[CH₂–CH(Ar)]–), which gives the exact mass of the polymer repeat unit.

\subsection{Conformers and pre-optimization}

For each motif, RDKit [RDKit] generates a set of 3D conformers, which are ranked with a force field. Each redox state is then pre-optimized in the gas phase with UMA (model uma-s-1p2p1) \cite{wood2025}, a machine-learned interatomic potential (MLIP) that accounts for charge and spin. UMA is used only to choose good starting structures (\emph{seeds}) for DFT; it never supplies a solvated energy. In this study up to three seeds per redox state were carried into DFT, and the lowest-energy result was kept. Three seeds were used for 262 of the 274 states of the 92 scored candidates (the 91-candidate pool plus one DMPZ added later; see Results); the code's default is one seed. The choice of seed matters: of the 268 states with more than one seed, the seeds end within 5 meV of each other in 154, but differ by more than 25 meV in 62 and by more than 0.1 eV in 11. An optional scan of the ring–ring torsion is available for flexible molecules.

\subsection{DFT geometry optimization}

Every redox state is \textbf{optimized in solvent}, using the SMD implicit-solvent model for acetonitrile (MeCN) \cite{marenich2009} and the geomeTRIC optimizer \cite{wang2016}, because a gas-phase structure can be wrong for charged or flexible species. Each seed is first pre-optimized with GFN2-xTB \cite{bannwarth2019}, a fast semiempirical tight-binding method, in implicit MeCN (ALPB, the analytical linearized Poisson–Boltzmann model \cite{ehlert2021}); the solvent also keeps dianions bound, which in the gas phase can lose an electron. Geometries are then optimized with DFT at r²SCAN-D4/def2-SVP (def2-SVPD for anions) in SMD \cite{furness2020,caldeweyher2019,weigend2005,rappoport2010}. The self-consistent-field (SCF) calculations use density fitting (RI-J) \cite{weigend2006}, a convergence threshold of 1 × 10⁻⁹, and up to 200 iterations; open-shell cases that fail to converge are retried with stronger convergence aids (level shifting and damping, then a second-order solver).

An open-shell calculation can settle on an excited electronic solution, depending on its starting guess for the orbitals. Each open-shell state is therefore started from more than one guess and the lowest converged solution is kept. For organic molecules the guesses are the standard superposition-of-atoms guess (minao) and a Hückel guess; if these disagree by more than 1 meV, or one fails to converge, an atomic-density guess is also run and the state is flagged for review. On 250 organic open-shell states this pair always reached the lowest solution found by any guess (SI S5). Metal-containing open shells use all three guesses, plus extra starting points built for the Fe doublet of ferrocenium. (The pool's energies were computed with all three guesses for every open shell. On the 250 test states the two-guess policy finds the same lowest solution; see SI S5.) The def2 effective core potential (ECP) \cite{weigend2005} is added automatically for elements with Z ≥ 37. The redox states of a motif are independent and can be optimized at the same time, one GPU per state.

\subsection{Production energies}

The production energy of every charge state, and of the ferrocene reference, is computed at one level: \textbf{ωB97M-V/def2-TZVPD} \cite{mardirossian2016,rappoport2010}, a functional that includes the VV10 nonlocal dispersion term \cite{vydrov2010}, in SMD acetonitrile. Using one diffuse basis for all charge states is deliberate: a mixed scheme (def2-TZVP for neutrals, def2-TZVPD for anions) would add diffuse functions to only one side of every reduction, so errors would not cancel in the potential. Each energy is stored with a fingerprint of everything that determines it (coordinates, charge, multiplicity, functional, basis, dispersion, solvent model, density fitting, and convergence settings). An energy left over from the optimization is reused only if its fingerprint matches these settings exactly; otherwise a new single point is computed. Both the solvent (SMD) and gas-phase energies are kept. Their difference is the solvation energy, dG\_solv = E\_SMD − E\_gas, and the gas-phase highest occupied orbital (HOMO) energy is used to flag anions that would not be bound in the gas phase. The gas-phase energy is needed only for the selected seed, so it is computed once, after the seed is chosen.

\subsection{Thermal corrections: two tiers}

Free energies are G = E\_SMD + G\_thermal, where E\_SMD is the production energy in solvent. Two levels of thermal correction are used. The \textbf{xTB tier}, used for the whole pool, is a GFN2-xTB (a fast semiempirical tight-binding method) rigid-rotor–harmonic-oscillator (RRHO) correction at 298.15 K, computed on the DFT+SMD geometry: G\_thermal = ZPE + H\_thermal − TS, with ZPE the zero-point energy. It is cheap enough to apply identically to every molecule and state, and much of it cancels in the energy \emph{differences} we report. On the experimental benchmark it is no less accurate than the second tier: re-scoring the 39 benchmark potentials with it, with all electronic energies unchanged, gives an MAE of 0.232 V instead of 0.247 V (Results). The second tier, the \textbf{DFT-frequency tier}, computes the vibrational frequencies with DFT in solvent (r²SCAN-D4/def2-SVP(D)/SMD). It accepts a structure only if it is a true energy minimum, adds a spin-degeneracy term, and treats low-frequency vibrations with a quasi-RRHO (qRRHO) model \cite{grimme2012}. It needs a GPU and much disk space, so it is used to \emph{confirm} leads rather than across the pool. If a state lacks either its production energy or a valid thermal correction, its couple is marked incomplete; electronic energies alone are never used instead. A molecule and the ferrocene reference must use the same thermal tier, or the couple is flagged.

\subsection{Derived properties (kept separate)}

Each property is computed and reported separately; they are never combined.

\begin{itemize}
\item \textbf{Redox potential vs Fc.} For a one-electron couple O + e⁻ → R, dG = G(R) − G(O), E\_abs = −dG, and E\_vs\_Fc = E\_abs − E\_ref. The ferrocene/ferrocenium (Fc/Fc⁺) reference E\_ref is computed at the same level, with the same free-energy definition, as the molecules, so errors common to both cancel (RT/F = 0.025693 V, Nernst slope 0.059160 V).
\item \textbf{Usable capacity.} Theoretical specific capacity C\_spec = nF/(3.6·M) [mAh g⁻¹], with M the neutral repeat-unit mass of the graft (above) and counterions excluded, the usual active-material convention. An electron counts as usable only if its potential lies inside the stability window of the electrolyte, tetrabutylammonium hexafluorophosphate (TBAPF₆) in MeCN ([−3.0, 2.0] V vs Fc); a candidate whose second electron falls outside the window is credited with one. Couples below 0 V vs Fc count as anolytes and those above as catholytes. A salt-loaded variant that adds the counterions (n PF₆⁻ for catholytes, n tetrabutylammonium TBA⁺ for anolytes) is reported in the SI. Because TBA⁺ (242 g mol⁻¹) is much heavier than PF₆⁻ (145 g mol⁻¹), that variant penalizes anolytes more than catholytes and must not be used to compare the pools; it does not change Pareto-front membership within either pool.
\item \textbf{Outer-sphere λ\_o.} A continuum estimate of the solvent's reorganization energy. The single-sphere Born/Marcus expression \cite{born1920,marcus1956}, with the radius of the sphere that has the molecule's solvent-accessible surface area (SASA), is reported for every candidate (SI S7). As a check, a nonequilibrium polarizable continuum model (IEF-PCM) \cite{cances1997} that uses the molecule's actual cavity shape was run for 10 couples of 5 pool molecules. It gives 1.03 to 1.39 times the Born value (ε\_s = 37.5, ε\_op = 1.806 in both). Values are per half-reaction. How far these continuum estimates are from explicit-solvent values is not tested here.
\item \textbf{Inner-sphere λ\_i.} Nelsen four-point scheme \cite{nelsen1987}: each redox state's energy at its own optimized geometry and at its partner's geometry. The values reported here use gas-phase energies at geometries optimized in solvent, from single conformers. Those geometries are not energy minima for gas-phase energies, so λ\_i can come out negative, and for anions the gas-phase state can be unbound. λ\_i is therefore treated as qualitative and used for no kinetic or reversibility claim. The current code evaluates all four points in solvent instead, which removes these problems; the pool's λ\_i values have not yet been recomputed that way.
\item \textbf{Disproportionation.} dG\_disp = G(R\textasciicircum{}\{q+1\}) + G(R\textasciicircum{}\{q−1\}) − 2G(R\textasciicircum{}q) = F·(E\_high − E\_low), i.e. the spacing between a molecule's two waves, which does not depend on the reference. It is reported as computed, without correction.
\end{itemize}

\subsection{Reliability gating and quality control}

Every redox state is checked automatically (\texttt{quality/readiness.py}): the actual charge, atom composition, connectivity, stereochemistry, and spin parity must match the declared structure; the geometry must come from the current protocol; and a valid production single point and thermal term must exist. Imaginary thermal modes are flagged for review. Spin states are checked with the expectation value ⟨S²⟩ (spin contamination of at most 0.1 is allowed) and, where two spin states lie close in energy, by recomputing the state at the other spin multiplicity with DFT and comparing. We also check that every redox state of a molecule is present, that the extra geometries needed for λ\_i exist, and that all states use the same thermal correction. States that fail are \textbf{held, not dropped}: they are marked incomplete and keep their reason. In the 91-candidate pool, 269 of the 271 states pass every check the ranking relies on (Figure 3). The two that do not are dianion states of two imide anolytes whose spin state has not yet been confirmed by the checks described above. Those two candidates are kept in the pool with a flag; neither is on a Pareto front, and the fronts are unchanged without them.

\begin{figure}[htbp]
\centering
\includegraphics[width=0.98\textwidth]{figures/fig3_qc_funnel.png}
\caption{\textbf{How the registry of molecules narrows to the scored pool through automated quality checks.} The funnel has four tiers: registered graphs → state-complete → QC-clean → scored candidates. A redox state is \emph{state-complete} when it has a converged geometry and a valid production single point. It is \emph{QC-clean} when it also passes every check the ranking relies on: the realized charge, atom composition, connectivity, stereochemistry, and spin parity match the declared structure; the geometry was made with the current protocol; and a valid thermal correction exists (\texttt{quality/readiness.py}, \texttt{quality/spins.py}, \texttt{quality/stability.py}). \textbf{States that fail a check are held, not dropped}: they stay in the records with their reason.}
\label{fig:3}
\end{figure}

\subsection{Multi-property ranking}

Candidates are ranked on \textbf{Pareto fronts, with the anolyte and catholyte pools kept separate}. An ambipolar molecule enters both pools with only its own side's electrons. The fronts in this paper use two objectives, oriented so that higher is better: voltage (−E for anolytes, +E for catholytes) and theoretical capacity. A candidate is on the front if no other candidate in its pool is at least as good on both and better on one. The benchmark's mean absolute error (±0.25 V) is drawn on the figures as the uncertainty of each potential, rather than built into the front. We also tested a rule that takes uncertainty into account. Under it, candidate A beats candidate B only if A is clearly better on one property and not clearly worse on the other, where "clearly" means by more than the combined uncertainty (each family's benchmark root-mean-square error, RMSE). Capacity has no computational uncertainty, so any candidate with higher capacity then beats every candidate whose potential lies within the combined uncertainty (about 0.2–0.5 V, depending on the families), and each front collapses to its single highest-capacity member. We therefore report the plain fronts and ask the reader to treat differences in potential smaller than the shown uncertainty as unresolved. The ranking code also supports disproportionation stability and −λ\_i as objectives; λ\_i is not used here because it is qualitative (above). Solubility and synthetic-accessibility proxies are carried as annotations, not objectives. There is deliberately \textbf{no weighted-sum score}: it would need arbitrary weights and would change whenever the candidate set changes. The front and the raw values are reported instead.

\subsection{Reproducibility and software}

Every stage is launched from a command registry: \texttt{motifs.commands()} gives the exact command line for each stage. Each run writes a \texttt{run\_manifest.json} that records the configuration and its hash, the selected molecules, the host, the GPU and CPU assignment, thread limits, the xTB program, the number of conformer seeds, the thermal tier, and the return code and timing of every step. A resumed run must use the same inputs: new molecules may be added, but a changed molecule stops the run. Reproducibility is tested in two ways. Five worked examples (\texttt{examples/}) must reproduce stored results byte for byte, and a test suite (780 passing) checks the code. The thermochemistry code, the xTB program, and the inputs of the worked examples are each checked against a stored SHA-256 hash before use, so any change to them is detected.

Each calculation ran on its own GPU, assigned through a reservation tool, with a fixed number of CPU threads.

The pipeline normally checks the arithmetic of its exported tables. The table used in this paper was assembled separately, so we checked it independently. Every couple of all 92 candidates was recomputed from the stored energies and geometries (\texttt{scripts/validation/score\_pool\_legacy\_tier.py}), and all agree with the table to within 0.5 mV, the size of its rounding (SI S7). The calculations use PySCF 2.14.0 \cite{sun2020} with GPU acceleration from gpu4pyscf 1.8.1 \cite{li2025,wu2025}, xtb 6.7.1, and RDKit 2026.03.5. The comparison of computed potentials with experiment is reported in the Results, together with the code version that produced it.

\begin{table}[htbp]
\caption{Level of theory and standard-state conventions}\label{tab:1}
\small
\begin{tabularx}{\textwidth}{@{}>{\raggedright\arraybackslash}X>{\raggedright\arraybackslash}X@{}}
\hline
Item & Setting \\ \hline
Conformer pre-\allowbreak{}optimization & RDKit conformers (force-\allowbreak{}field-\allowbreak{}ranked) → UMA MLIP, gas phase \\
Geometry optimization & r²SCAN-\allowbreak{}D4/def2-\allowbreak{}SVP (anions def2-\allowbreak{}SVPD), SMD(MeCN), geomeTRIC \\
Solvated pre-\allowbreak{}optimization & GFN2-\allowbreak{}xTB, ALPB(MeCN) warm start \\
Production energy & ωB97M-\allowbreak{}V (with VV10 dispersion)/def2-\allowbreak{}TZVPD, \textbf{same diffuse basis for all charge states and Fc} \\
Solvation model & SMD, acetonitrile (ε\_\allowbreak{}s = 37.5, ε\_\allowbreak{}op = 1.806, n = 1.344) \\
SCF & RI-\allowbreak{}J density fitting; convergence 1 × 10⁻⁹; up to 200 iterations; open shells from two or more guesses, lowest kept \\
Heavy-\allowbreak{}element treatment & def2 ECP for Z ≥ 37 \\
Thermal correction, xTB tier (whole pool) & GFN2-\allowbreak{}xTB RRHO on the DFT+SMD geometry, 298.15 K \\
Thermal correction, DFT-\allowbreak{}frequency tier (lead confirmation; benchmark) & DFT frequencies in solvent (r²SCAN-\allowbreak{}D4/def2-\allowbreak{}SVP(D)/SMD) with quasi-\allowbreak{}RRHO \\
Free energy & G = E\_\allowbreak{}SMD + G\_\allowbreak{}thermal (one matched thermal definition per molecule and Fc) \\
Temperature & 298.15 K (RT/F = 0.025693 V; Nernst slope 0.059160 V) \\
Potential reference & Fc/Fc⁺, computed at the same level and thermal tier as the molecules \\
Electrolyte / window & TBAPF₆ / MeCN; [−3.0, 2.0] V vs Fc; anolyte/catholyte divider 0 V \\
Outer-\allowbreak{}sphere λ\_\allowbreak{}o & Born (SASA radius) for every candidate; molecular-\allowbreak{}cavity IEF-\allowbreak{}PCM check on 5 molecules (10 couples) \\
Inner-\allowbreak{}sphere λ\_\allowbreak{}i & Nelsen four-\allowbreak{}point; reported values use gas-\allowbreak{}phase energies on SMD geometries (qualitative) \\
Ranking & Pareto fronts on voltage and capacity, anolyte/catholyte pools separate, uncertainty shown; \textbf{no weighted-\allowbreak{}sum score} \\
\hline
\end{tabularx}
\end{table}

\emph{Solvent constants are those recorded in \texttt{config/project.json}: static permittivity ε\_s = 37.5 and refractive index n\_D = 1.344, which give the optical permittivity ε\_op = n\_D² = 1.806. The Pekar factor 1/ε\_op − 1/ε\_s enters the outer-sphere λ\_o. The value n\_D = 1.344 is the standard tabulated value for MeCN at 20 °C. It changes little with temperature (dn/dT ≈ −4 × 10⁻⁴ K⁻¹), so the \textasciitilde{}5 K difference from the 298.15 K thermochemistry shifts n\_D by ≈ 0.002 and ε\_op by < 0.01, a negligible effect on λ\_o. We use the 20 °C n\_D with the 298.15 K free energies, the usual practice.}

\section{Results and Discussion}

\subsection{Validation against experiment}

The credibility of any ranking depends on whether the pipeline reproduces measured electrochemistry. Before applying it to the candidate pool, we therefore test it against literature redox potentials measured under matching conditions. The validation is deliberately retrospective and uncorrected: we compare the pipeline's computed potentials of \emph{free parent} molecules against experimental values measured in acetonitrile and referenced to ferrocene (Fc), the reference recommended for nonaqueous solvents \cite{gritzner1984}. The molecules go through the same build, DFT+SMD optimization, and single-point steps as the screening candidates. The thermal correction differs: the benchmark as first computed uses the DFT-frequency tier, while the pool uses the cheaper xTB tier, so we report the benchmark at both (below). The experimental panel is hand-curated from seven studies \cite{huynh2016,wang2020,cabral2023,kim2022,daub2021,attanayake2019,gerken2015}. Each measured wave carries its source, DOI, table or page, and original reference electrode. Whether a value is used depends only on what the source reports, following rules written into the benchmark configuration (\texttt{config/benchmark\_mecn.py}). Only (quasi)reversible half-wave potentials reported against Fc, directly or through the authors' own calibration, are scored (category A). Half-peak potentials, electrolytes with added water, unverified conversion constants, waves seen only as a cathodic peak, and compounds with unspecified stereocenters or uncertain identity are set aside with the reason recorded. Every category-A value (an \emph{endpoint}: one measured wave of one molecule) from these sources is included, no molecule is excluded, and all outliers are kept. Potentials are computed as E\_vs\_Fc = E\_abs − E\_ref (Methods), with the ferrocene reference computed at the same level of theory and with the same free-energy definition (G = E\_SMD + G\_thermal) as the molecules. No correction or fitted offset is applied. The panel was not held out from method development, so it is a description of accuracy, not a test of a fitted model.

\textbf{Headline result.} Across the fixed category-A panel the pipeline reproduces experimental potentials with a mean absolute error of \textbf{0.247 V} and a root-mean-square error (RMSE) of 0.279 V. The \textbf{systematic bias is +0.19 V} (computed potentials sit above experiment), and the largest error is 0.547 V. The panel comprises \textbf{39 scored endpoints (redox couples) drawn from 30 molecules and seven independent studies}. Averaging first within each molecule and then across molecules gives an MAE of 0.263 V. The bias depends on the direction of the couple: the 33 reduction endpoints (anolyte-like) give MAE 0.239 V and bias +0.17 V, whereas the six oxidation endpoints (five phenothiazine, one nitroxide; catholyte-like) give MAE 0.289 V and bias +0.29 V. Uncertainty was estimated by bootstrap resampling of whole molecules, so that the waves of one molecule stay together (2000 replicates, seed 20261021): the 95\% interval is [0.213, 0.286] V for the MAE and [0.128, 0.257] V for the bias. Resampling whole studies instead, which gives wider intervals because there are only seven, yields [0.178, 0.306] V and [0.053, 0.285] V (SI). These are descriptive intervals for the observed panel, not prediction intervals for a new molecule.

\textbf{At the pool's thermal tier.} We re-scored the same 39 potentials with the xTB thermal correction used for the pool, for the molecules and for ferrocene alike, keeping every electronic energy unchanged. This gives an MAE of 0.232 V, an RMSE of 0.266 V, and a bias of +0.19 V (reductions: MAE 0.224 V, bias +0.18 V; oxidations: MAE 0.278 V, bias +0.28 V). Individual potentials move by −0.06 to +0.11 V between the two tiers, and by +0.003 V on average. The cheaper tier used for the pool is therefore no less accurate on this panel. The re-scoring script and its record are in the SI (\texttt{scripts/validation/benchmark\_tier\_comparison.py}); recomputing the DFT-frequency tier the same way reproduces every published potential exactly.

\textbf{Interpretation.} The dominant error is a \emph{systematic} positive offset, not scatter: the bias (+0.19 V) is a large fraction of the MAE. Splitting the errors by wave shows where it sits. The first reduction or oxidation (wave 1) carries nearly all of it (MAE 0.266 V, bias +0.266 V). The second wave (wave 2) has a smaller error with almost no net offset (MAE 0.185 V, bias −0.064 V). These wave statistics are for the published benchmark (DFT-frequency tier). A shared offset does not change the order of candidates, so comparisons within a pool are more reliable than absolute values, and that is how the screen below is used. We do not attribute the offset to a single cause. A part of it could come from the ferrocene reference, which shifts every potential equally. But the spacing between the two waves of a molecule (E1 − E2) does not depend on the reference, and it is computed larger than measured in eight of the nine molecules with two measured waves. A reference error alone therefore cannot explain the residuals. Errors in how solvation and ion pairing are treated for each charge state probably contribute.

Three limits apply. The panel tests \emph{free parent} molecules, not grafts or polymers. It contains no naphthalene diimide (NDI), because the only NDI entry was excluded for uncertain identity and unspecified stereocenters. And we cannot yet say how much of the bias comes from the ferrocene reference and how much from solvation.

The panel contains no dihydrophenazine, so it does not by itself calibrate the catholyte family used in the showcase below. As a family-specific check we computed the exact measured parent, 5,10-bis(2-methoxyethyl)-5,10-dihydrophenazine (BMEPZ; E₁ = −0.29 V, E₂ = +0.50 V vs Fc in MeCN \cite{kwon2019}), with the xTB thermal tier used for the pool. The computed couples are −0.06 and +0.87 V, errors of +0.23 and +0.37 V (MAE 0.30 V, bias +0.30 V). This is one molecule, so it is a check rather than a benchmark. It agrees with the oxidation side of the panel at the same xTB tier, where oxidations are too positive by +0.28 V on average. The panel's only second oxidation (of N-ethyl-3,7-dimethoxyphenothiazine) is too positive by 0.48 V at that tier, so the larger BMEPZ error on its second wave is not unexpected. The BMEPZ check therefore implies that the absolute DMPZ potentials reported below are likely about 0.3 V too positive. The full record (per-state energies, thermal corrections, optimized geometries, and the ferrocene reference with the hashes of its source records) is provided in the SI and regenerated by \texttt{scripts/validation/bmepz\_anchor.py}.

\begin{figure}[htbp]
\centering
\includegraphics[width=0.49\textwidth]{figures/fig4_validation_parity.png}
\includegraphics[width=0.49\textwidth]{figures/fig4b_residuals.png}
\caption{Computed versus experimental redox potentials (V vs Fc/Fc+, in MeCN) for the tier-A benchmark panel. Each point is one redox endpoint (a single measured wave / couple), coloured by redox family (quinone, aromatic imide, phenothiazine, viologen, nitroxide). The solid line is the 1:1 identity; the shaded envelope is the +/-MAE agreement band. Over \textbf{n = 39 endpoints from 30 molecules and 7 independent studies}, the protocol reaches \textbf{MAE = 0.247 V, RMSE = 0.279 V}, with a \textbf{systematic bias of +0.19 V} (computed potentials run high); the largest single error is 0.55 V (the second oxidation of N-ethyl-3,7-dimethoxyphenothiazine), kept in the panel. By family: quinone n = 22, MAE 0.286 V; aromatic imide n = 9, MAE 0.158 V; phenothiazine n = 5, MAE 0.296 V; viologen n = 2; nitroxide n = 1. Level of theory: geometries at r2SCAN-D4/def2-SVP(D)/SMD(MeCN), energies at wB97M-V/def2-TZVPD/SMD(MeCN).}
\label{fig:4}
\end{figure}

\textbf{Which code produced these numbers.} The benchmark was computed with the version of the code at commit \texttt{46b39c3b} (22 October 2026). That commit is preserved in the public repository under the tag \texttt{legacy/pre-refactor-main}, and the benchmark record, which records hashes of its inputs, is provided in the SI. The released code is a reorganized version of that code. Its five worked examples reproduce stored results byte for byte. Two later changes to the DFT step were tested directly: the gas-phase energy is now computed only for the selected conformer, and radicals use the new guess policy. Both reproduce the previous energies to within 10⁻¹² hartree (Ha). We have not, however, rerun the benchmark's DFT calculations with the released code; the tagged version is the one to use to reproduce the benchmark exactly.

\subsection{Mapping the 91 candidates: separate anolyte and catholyte fronts}

For each of the 91 grafted candidates the pipeline returns a quality-checked record of properties: the potential versus ferrocene (Fc) in TBAPF₆/MeCN, the usable capacity, outer- and inner-sphere reorganization estimates, and disproportionation free energies. From these the design space can be read directly (Figure 5). The pool is strongly anolyte-dominated: 85 of the 91 candidates reduce (n-type anolytes) and only 6 oxidize into an accessible catholyte window (p-type). By family the set is 36 imides, 29 quinones, 18 viologens (bipyridiniums), 5 dihydrophenazines, 2 mono-imides, and 1 phenothiazine; 89 of the 91 are two-electron couples, the two exceptions being the single-wave phthalimide and 1,8-naphthalimide controls. The candidates combine 23 motifs derived from structures in the D3TaLES database \cite{duke2022}, grafted versions of literature redox cores, and designed variants that change one substituent, linker, or cap at a time. The composition reflects the source chemistry: graftable, well-precedented n-type cores (naphthalene- and pyromellitic-diimides, benzo-, naphtho-, and anthraquinones) are abundant, whereas graftable p-type cores are scarce.

\begin{figure}[htbp]
\centering
\includegraphics[width=0.49\textwidth]{figures/fig5a_anolyte_pareto.png}
\includegraphics[width=0.49\textwidth]{figures/fig5b_catholyte_pareto.png}
\caption{\textbf{Redox potential versus usable capacity for the 91-candidate grafted-motif pool, resolved into the two separate electrolyte roles.} Each point is one candidate, placed by its two computed performance quantities kept on \emph{independent} axes — reduction/oxidation potential \emph{E} vs Fc/Fc+ (x) and theoretical gravimetric capacity (y; nF/3.6M on the neutral repeat-unit mass, counterions excluded — the salt-loaded variant, which leaves front membership unchanged, is in the SI) — and coloured by redox family. \textbf{No weighted-sum or single-number figure of merit is used:} potential and capacity are never combined, so the trade-off is read directly from the plane rather than from a single ranked score. The Pareto-optimal front within each pool (step line, outlined numbered markers keyed to real chemical names) is the set of candidates not dominated on both axes simultaneously — for anolytes, more negative potential \emph{and} higher capacity; for catholytes, more positive potential \emph{and} higher capacity. Horizontal whiskers on front members span ±0.25 V (MAE 0.247), the validation-benchmark potential error (MAE over 39 tier-A redox endpoints); front membership that would flip under a potential shift smaller than this is not physically meaningful, which the whiskers make explicit without collapsing the two axes. \textbf{(a) Anolyte pool (85 n-type candidates).} Quinones (AQ/NQ/BQ), imides (PMDI/NDI) and multi-electron viologens densely populate the low-potential region; the front is a smooth potential/capacity trade-off anchored by an O-grafted benzoquinone at the high-capacity end and a pyromellitic-diimide (PMDI) graft at the most-negative-potential end. The two 1-electron mono-imide controls (open circles) are shown for reference and excluded from the front. \textbf{(b) Catholyte pool (6 p-type candidates).} The p-type space — 5,10-dihydrophenazines (DMPZ) plus a single dimethoxy-phenothiazine (PTZ) — is strikingly sparse: six points against the anolytes' eighty-five, clustered at modest potential with a wide, essentially unexplored high-potential region (shaded). This asymmetry, not any score, is the figure's message: the design space is anolyte-dominated and the catholyte side is the open gap.}
\label{fig:5}
\end{figure}

Candidates are ranked not by a single score but on Pareto fronts computed \textbf{separately for the anolyte and catholyte pools}. A redox potential is not "more is better" in general: an anolyte is better the more negative its potential, a catholyte the more positive. Within each pool the two objectives are voltage and theoretical capacity (mAh g⁻¹, active mass), and a candidate is on the front if no other candidate is at least as good on both and better on one (Methods). The ±0.25 V uncertainty is drawn on each front member; differences in potential smaller than that should be read as unresolved. We deliberately avoid any weighted-sum score, which would need arbitrary weights and would change whenever the candidate set changes. The front and the raw values are given instead (Figure 5; full table in the SI). Inner-sphere reorganization (λᵢ) is carried only qualitatively and is not used in the ranking, for the reasons given in the Methods.

Figure 5 presents the two Pareto panels side by side. The anolyte panel (85 candidates) is densely populated across a wide range of potentials. Imides span −0.81 to −1.54 V vs Fc and quinones −0.83 to −1.49 V. Pyromellitic diimides (PMDI; e.g. the ether-tethered PMDI at −1.54 V) and alkoxy-anthraquinones (≈ −1.49 V) sit at the most negative edge, and the viologens lie between −0.63 and −0.99 V. Theoretical capacities of the two-electron anolytes span ≈ 90–200 mAh g⁻¹, highest for the light benzo- and naphthoquinones and falling as solubilizing oligoether or long-alkyl caps add inert mass. The catholyte panel, by contrast, is almost empty: six points, five of them dihydrophenazines between +0.30 and +0.62 V and one 3,7-dimethoxyphenothiazine at +0.78 V. This visual sparsity — a crowded anolyte front against a nearly bare catholyte front — is the central result of the screen, and it defines the \textbf{catholyte gap} that the next step is designed to fill.

\subsection{Filling the catholyte gap with grafted dihydrophenazines}

To populate the empty catholyte front we grafted a single scaffold, 5,10-dihydrophenazine (DMPZ), and varied one substituent at a time, so that each change in potential can be traced to one substituent rather than to several differences at once (Figure 6). Every member carries the same 4-methylbenzyl (Merrifield proxy) cap on one nitrogen and the same two-electron sequence, neutral → +1 → +2, isolating the effect of the N′-substituent and of ring substitution. Table 2 lists the series.

\begin{figure}[htbp]
\centering
\includegraphics[width=0.49\textwidth]{figures/fig6_dmpz_design_rule.png}
\includegraphics[width=0.49\textwidth]{figures/fig6b_dmpz_catholyte_space.png}
\caption{\textbf{A matched single-substituent rule for tuning a grafted 5,10-dihydrophenazine (DMPZ) catholyte.} Mean redox potential versus ferrocene (E vs Fc/Fc⁺) in TBAPF₆/MeCN for a family built by grafting one scaffold — DMPZ with a 4-methylbenzyl (Merrifield-proxy) cap on one nitrogen and a two-electron neutral → +1 → +2 sequence of states — and varying a single substituent at a time, so that each shift reads as a matched delta rather than a confounded whole-molecule comparison. Potentials (E = mean of the two one-electron couples) are read live from \texttt{results/validation/full\_pool\_scored\_legacyG\_MeCN\_TBAPF6.csv}; no value is altered. DMPZ members are drawn in orange and the distinct 3,7-dimethoxyphenothiazine (PTZ) in purple (fixed family colours across the figure set). Members are ordered by potential; the dashed line marks the N′-ethyl DMPZ reference (+0.43 V), and the parenthetical value in each annotation is the matched shift relative to that reference (red where a substituent raises E, green where it lowers E).}
\label{fig:6}
\end{figure}

\begin{table}[htbp]
\caption{Grafted DMPZ series and the 3,7-dimethoxyphenothiazine. E vs Fc is the mean of the two one-electron couples (xTB thermal tier); capacity uses the neutral repeat-unit mass. Identifiers are those used in the records and in SI Table S7.}\label{tab:2}
\footnotesize
\begin{tabularx}{\textwidth}{@{}>{\raggedright\arraybackslash}X>{\raggedright\arraybackslash}Xr>{\raggedright\arraybackslash}X>{\raggedleft\arraybackslash}X@{}}
\hline
Candidate & N′ / ring substitution & E vs Fc (V) & Couples (V) & Capacity (mAh g⁻¹) \\ \hline
\texttt{exp\_\allowbreak{}dmpz\_\allowbreak{}mb\_\allowbreak{}ethyl} & N′-ethyl & +0.43 & −0.03 / +0.89 & 164.2 \\
\texttt{exp\_\allowbreak{}dmpz\_\allowbreak{}mb\_\allowbreak{}methyl} & N′-methyl & +0.47 & +0.03 / +0.91 & 171.6 \\
\texttt{hyp\_\allowbreak{}dmpz\_\allowbreak{}mb\_\allowbreak{}methoxyethyl} & N′-(2-\allowbreak{}methoxyethyl) & +0.47 & +0.01 / +0.94 & 150.4 \\
\texttt{hyp\_\allowbreak{}dmpz\_\allowbreak{}mb\_\allowbreak{}trifluoroethyl} & N′-(2,\allowbreak{}2,\allowbreak{}2-\allowbreak{}trifluoroethyl) & +0.62 & +0.15 / +1.10 & 140.9 \\
\texttt{hyp\_\allowbreak{}dmpz\_\allowbreak{}27dimethyl\_\allowbreak{}ethyl} & 2,\allowbreak{}7-\allowbreak{}dimethyl, N′-ethyl & +0.30 & −0.14 / +0.73 & 151.2 \\
\texttt{hyp\_\allowbreak{}dmpz\_\allowbreak{}27dimethoxy\_\allowbreak{}ethyl} & 2,\allowbreak{}7-\allowbreak{}dimethoxy, N′-ethyl & +0.21 & −0.22 / +0.64 & 138.7 \\
\texttt{exp\_\allowbreak{}ptz\_\allowbreak{}37dimethoxy} & 3,\allowbreak{}7-\allowbreak{}dimethoxyphenothiazine & +0.78 & +0.37 / +1.19 & 142.7 \\
\hline
\end{tabularx}
\end{table}

Read as shifts relative to the N′-ethyl DMPZ reference (+0.43 V), the series gives a consistent \textbf{design rule}. An electron-withdrawing N′-substituent raises the potential: replacing ethyl with 2,2,2-trifluoroethyl shifts E by +0.20 V (to +0.62 V), while N′-methyl and N′-(2-methoxyethyl) barely move it (+0.04 to +0.05 V). Conversely, electron-donating ring substitution lowers it: 2,7-dimethyl brings E down by ≈ 0.13 V (to +0.30 V) and 2,7-dimethoxy by ≈ 0.22 V (to +0.21 V) relative to the same N′-ethyl parent. The two handles push the potential in opposite directions, and together they span ≈ +0.21 to +0.62 V across the grafted series; the distinct 3,7-dimethoxyphenothiazine scaffold reaches +0.78 V. Whether the two effects add when combined in one molecule was not tested, since no member of the series carries both. These light two-electron phenazine cores give theoretical capacities of ≈ 140–170 mAh g⁻¹. That is comparable to the anolyte pool, so the catholyte side is not a capacity bottleneck.

Two cautions apply to these numbers. First, the matched shifts are the robust quantity: within one scaffold, errors common to the family cancel. Second, the absolute potentials are probably too high. The BMEPZ check above puts the DMPZ couples about 0.3 V too positive, so the series likely sits nearer −0.1 to +0.3 V vs Fc in experiment. Re-plotting the catholyte Pareto panel with these points shows the gap being filled by a structured family rather than by isolated outliers (Figure 6).

Like every value in the pool, these potentials use the ranking thermal tier. The pipeline's design calls for confirming leads with the more expensive DFT-frequency tier, and that has not yet been run for any DMPZ member. On the benchmark the two tiers differ by −0.06 to +0.11 V per potential (above), which is a guide to how much these values could move. The 91-candidate pool behind Figure 5 contains six catholytes (five DMPZ and one phenothiazine). The 2,7-dimethoxy DMPZ (E ≈ +0.21 V) was added after that pool was fixed, so the scored table has 92 rows. Figure 5 shows the 91-candidate pool; Figures 6 and 7 also include this seventh catholyte.

\subsection{Proposed non-aqueous full cells}

Because every potential in both pools is on the same Fc scale, candidates can be paired directly, and a nominal cell voltage read off as E(catholyte) − E(anolyte) (Figure 7). Table 3 gives four example pairings, from the mean potentials in the scored table.

\begin{figure}[htbp]
\centering
\includegraphics[width=0.98\textwidth]{figures/fig7_cell_voltage.png}
\caption{\textbf{Proposed non-aqueous full cells read off a shared ferrocene scale (computed proposals, not demonstrated devices).} Because every candidate potential in both screening pools is referenced to the same ferrocene (Fc/Fc⁺) scale, candidates can be paired directly and a nominal cell voltage read off as \emph{E}₍cell₎ = \emph{E}(catholyte) − \emph{E}(anolyte). The energy axis is potential \emph{E} (V vs Fc) in TBAPF₆/MeCN. The upper green band is the catholyte pool (p-type): the grafted 5,10-dihydrophenazines (DMPZ) and one 3,7-dimethoxyphenothiazine (PTZ), spanning +0.21 to +0.78 V, with the viable DMPZ window near +0.3 to +0.6 V. The lower blue band is the anolyte pool (n-type quinone/imide/pyridine; 85 candidates shown as a rug), spanning −1.63 to −0.63 V. Three representative pairings are drawn as vertical arrows annotated with the resulting cell voltage: N′-ethyl DMPZ (+0.43 V) × alkoxy-anthraquinone [2-(4-methylbenzyloxy)-6-methoxyanthracene-9,10-dione, −1.49 V] → \textbf{1.91 V}; N′-ethyl DMPZ (+0.43 V) × ether-tethered pyromellitic diimide [PMDI ether\_C2, −1.54 V] → \textbf{1.97 V}; and the more electron-poor N′-(2,2,2-trifluoroethyl) DMPZ (+0.62 V) × the same ether-PMDI (−1.54 V) → \textbf{2.17 V} (marked with a star), extending the proposal beyond 2.1 V. These are computed values. Subtracting each family's measured bias at the pool's thermal tier (DMPZ +0.30 V from the BMEPZ check; quinones +0.24 V and imides +0.05 V from the benchmark) gives 1.85, 1.73 and 1.92 V; the PMDI pairs drop most because the imide bias is small. Potentials are the mean of the two one-electron couples from the scored pool (\texttt{results/validation/full\_pool\_scored\_legacyG\_MeCN\_TBAPF6.csv}, xTB-RRHO tier); each annotated voltage is the plotted arithmetic difference \emph{E}(catholyte) − \emph{E}(anolyte). \textbf{These are COMPUTED PROPOSALS derived from two computed half-cell potentials — not demonstrated cell voltage, attainable capacity, cycling, Coulombic efficiency, membrane crossover of the small grafted proxies, charged-state stability, or synthetic feasibility of the exact grafted structures; a capped graft is not its soluble parent or an oligomer.} Because catholyte and anolyte potentials move independently (the substituent design rule of Figure 6), the cell voltage is a tunable target rather than a single number.}
\label{fig:7}
\end{figure}

\begin{table}[htbp]
\caption{Proposed full cells. The bias correction subtracts each family's measured error at the pool's thermal tier (see text).}\label{tab:3}
\small
\begin{tabularx}{\textwidth}{@{}>{\raggedright\arraybackslash}X>{\raggedright\arraybackslash}Xrr@{}}
\hline
Catholyte & Anolyte & Cell (V) & Corrected (V) \\ \hline
N′-ethyl DMPZ (+0.43 V) & 2-\allowbreak{}(4-\allowbreak{}methylbenzyloxy)-\allowbreak{}6-\allowbreak{}methoxyanthraquinone (−1.49 V) & 1.91 & 1.85 \\
N′-methyl DMPZ (+0.47 V) & 1-\allowbreak{}(4-\allowbreak{}methylbenzyloxy)anthraquinone (−1.48 V) & 1.95 & 1.89 \\
N′-ethyl DMPZ (+0.43 V) & ether-\allowbreak{}tethered pyromellitic diimide (−1.54 V) & 1.97 & 1.73 \\
N′-(2,\allowbreak{}2,\allowbreak{}2-\allowbreak{}trifluoroethyl) DMPZ (+0.62 V) & ether-\allowbreak{}tethered pyromellitic diimide (−1.54 V) & 2.17 & 1.92 \\
\hline
\end{tabularx}
\end{table}

Computed potentials are too positive on both sides of the cell, so the two errors partly cancel: the error in a cell voltage is the catholyte family's bias minus the anolyte family's bias. The correction therefore subtracts each family's measured bias at the pool's thermal tier: +0.30 V for DMPZ (the BMEPZ check), +0.24 V for quinones, and +0.05 V for imides (benchmark). For the anthraquinone pairs the two biases nearly cancel, so the cell voltage barely changes. For the pyromellitic-diimide pairs they do not: the imide bias is small, so the correction lowers those cells by about 0.25 V. Taken together, the proposed cells lie at about 1.7–1.9 V. The design rule above makes the voltage adjustable: catholyte and anolyte potentials can be shifted independently.

We state the limitation plainly. A cell voltage computed from two half-cell potentials is a \textbf{proposal}, not a demonstrated device. It says nothing about capacity under cycling, Coulombic efficiency, or membrane crossover of the small grafted molecules. Nor does it show that the charged states stay stable over many cycles, or that the exact grafted structure can be made; each candidate's record lists the open questions about its synthesis. A capped graft is also not the same as its soluble parent or an oligomer. What the screen does provide is a quality-checked, internally consistent map of where these chemistries sit, a substituent rule for moving a DMPZ catholyte toward a target potential, and a concrete set of anolyte/catholyte pairings that experiments can test.

\subsection{Computational cost and throughput}

Every pipeline run records the start and end time of each stage, and the DFT stage additionally records the time spent in each conformer seed's pre-optimization, geometry optimization, solvated and gas-phase single points, and thermal correction. Table 4 summarizes the measured cost of five of the grafted DMPZ catholytes (43–52 atoms; 1031–1196 basis functions at def2-TZVPD). Each was taken from its SMILES string (a text encoding of the structure) to a ranked property record, with three redox states (neutral, radical cation, dication) and three conformer seeds per state. The per-run records are in \texttt{manuscript/tables/throughput\_measured.csv}, regenerated by \texttt{scripts/validation/throughput\_report.py}.

\begin{table}[htbp]
\caption{Measured per-candidate cost (GPU-minutes; one task per GPU, two CPU threads per worker). DFT is the sum over the three redox states; "λᵢ cross points" are the extra single points at the partner geometry needed for λ\_i.}\label{tab:4}
\footnotesize
\begin{tabularx}{\textwidth}{@{}>{\raggedright\arraybackslash}X>{\raggedright\arraybackslash}Xr>{\raggedleft\arraybackslash}X>{\raggedleft\arraybackslash}Xr@{}}
\hline
Candidate (atoms) & GPU & UMA & DFT (neu / +1 / +2) & λᵢ cross points & Total GPU-\allowbreak{}h \\ \hline
N′-methyl DMPZ (43) & V100 & 5.7 & 94.7 (19.8 / 47.1 / 27.8) & 9.7 & 1.8 \\
N′-trifluoroethyl DMPZ (46) & V100 & 6.0 & 126.1 (26.9 / 64.2 / 34.9) & 14.3 & 2.4 \\
2,\allowbreak{}7-\allowbreak{}dimethyl DMPZ (52) & V100 & 4.9 & 139.4 (34.4 / 65.4 / 39.6) & 14.5 & 2.6 \\
N′-ethyl DMPZ (46) & H100 & — & 47.9 (10.4 / 26.2 / 11.2) & 5.2 & 0.9 \\
N′-(2-\allowbreak{}methoxyethyl) DMPZ (50) & H100 & 4.9 & 83.8 (17.8 / 39.7 / 26.4) & 6.5 & 1.6 \\
\hline
\end{tabularx}
\end{table}

\emph{The N′-ethyl UMA stage was reused from an earlier run and is not counted. The 2,7-dimethoxy member's DFT stage was resumed across runs and has no single uninterrupted timing, so it is omitted.}

Three features of the cost structure follow from these measurements. First, DFT dominates: on the V100 runs it accounts for 87\% of GPU time, inner-sphere cross points for 9\%, and the MLIP pre-optimization for 4\%. Within DFT, the r²SCAN-D4/def2-SVP geometry optimizations take 57\% (V100) to 66\% (H100) of the time and the ωB97M-V/def2-TZVPD single points most of the rest. Second, the open-shell radical cation is the most expensive state, about half of all DFT time, because in these runs each open-shell single point was converged from three initial guesses in both solvent and gas phase. Third, cost rises steeply with size even across this narrow range: from 43 to 52 atoms (1031 to 1196 basis functions) the V100 DFT time grows by 1.47×, an apparent exponent of about 2.6 in the number of basis functions. With three molecules this exponent is indicative only. The H100 runs are 1.7–2× faster than the V100 runs for molecules of similar size, less than the difference in double-precision throughput between these GPUs. This suggests that part of each stage is limited by the two CPU threads per worker rather than by the GPU.

On this hardware a two-electron candidate costs about 0.9–2.6 GPU-hours, and the wall time per candidate is set by its slowest redox state (about 26–65 min) when each state has its own GPU. For example, the 91-candidate pool corresponds to roughly 80–240 GPU-hours at these rates. The pool's own records point to three reductions that leave every reported number unchanged. The first is the SCF guess policy described in the Methods: organic radicals now use two initial guesses instead of three, with automatic escalation when they disagree. On the 250 test states this loses no ground state, and it removes a third of the single-point SCFs for every open-shell state. Second, only the winning seed's gas-phase energy is used (for the solvation free energy and gas-phase diagnostics), so the pipeline now computes the gas single point once, on the selected seed, instead of for every seed. The resulting energies agree with the previous per-seed route to within 10⁻¹² Ha, the level of run-to-run floating-point noise. In the measured DMPZ runs this removes an estimated 13–15\% of DFT time. Third, the inner-sphere cross points, which feed a property reported only qualitatively, need be computed only for candidates that reach a Pareto front. Together these would remove an estimated 20–35\% of per-candidate GPU time; this is an estimate from the measured component times, not yet a measurement. Reducing the number of conformer seeds is not free. Across all 292 multi-seed states in the stored records (not only the pool), using only the first seed would have left 37 states more than 25 meV above the best seed's energy, and in 8 states the first seed failed outright.

We also tested, on the existing records, whether the much cheaper machine-learned potential (UMA) could replace DFT for ranking. For each type of couple (first or second reduction, first or second oxidation) we fitted a straight line from UMA gas-phase energy differences to the DFT potentials. Predicting each couple from a fit that leaves it out gives an MAE of 0.13 V against DFT over 174 couples (89 molecules). The rank correlation is 0.94 for first oxidations but only 0.63–0.72 for the anolyte reductions, whose DFT potentials span a narrow range (standard deviation 0.23 V). We then used these predictions as a pre-screen on the 60 two-electron anolytes that have both couples in the records (of the 83 two-electron anolytes in the pool). Every candidate predicted to lie within a set margin of the front was sent on to DFT. All seven anolyte front members of Figure 5 survive only when the margin is 0.3 V, and at that margin 43 of the 60 candidates still need DFT. In this pool, therefore, DFT remains necessary for the final ranking. An MLIP pre-screen becomes worthwhile for larger and chemically broader libraries, in which most candidates lie far from the front.

\section{Conclusions}

We have presented a reproducible computational pipeline for screening redox molecules to be grafted onto Merrifield-type polymers, and applied it to 91 candidates for non-aqueous organic redox flow batteries. The contribution is the method rather than any one molecule. The pipeline models each molecule as it is attached, with a 4-methylbenzyl cap standing in for the polymer. It computes every redox state at one level of theory (ωB97M-V/def2-TZVPD with SMD in acetonitrile) and reports potentials against ferrocene, together with capacity and reorganization energies. Before ranking, it checks every redox state for structure, spin, protocol, and completeness: 269 of the pool's 271 states pass, and the two that do not are flagged and affect no front. Each property is reported separately, and candidates are ranked on Pareto fronts of voltage and capacity, anolytes and catholytes apart, with the potential uncertainty shown and no weighted-sum score.

The applied results rest on a benchmark without any fitting. Across 39 endpoints (30 molecules, 7 studies) the computed potentials have a mean absolute error of 0.247 V and a systematic bias of +0.19 V (0.232 V with the cheaper thermal tier used for the pool). Because the error is largely a shared offset, the pipeline ranks candidates more reliably than it places absolute potentials. The 91-candidate screen shows a strongly anolyte-dominated space (85 anolyte / 6 catholyte; 89 of 91 are two-electron couples). Grafted 5,10-dihydrophenazine two-electron catholytes fill part of the sparse catholyte region. A series that changes one substituent at a time gives a simple design rule: electron-withdrawing N′-substituents raise the potential (2,2,2-trifluoroethyl to a computed +0.62 V vs Fc), while electron-donating ring groups lower it (2,7-dimethyl to +0.30 V, 2,7-dimethoxy to +0.21 V). A check against the measured parent compound suggests these absolute potentials are about 0.3 V too high. Pairing these catholytes with quinone and imide anolytes from the pool gives proposed cells of about 1.7–1.9 V after correcting for each family's measured bias.

These results are bounded by explicit limitations. The proposed full cells are proposals computed as the difference of two computed potentials; they are not demonstrated capacity, cycling, crossover, or synthetic feasibility, and a computed cell voltage does not establish any of these. A capped graft, a soluble parent, and an oligomer are three distinct models, and the capped graft is not the material. Inner-sphere reorganization (λ\_i) is reported only qualitatively: the present values rest on gas-phase energies evaluated on SMD geometries from single conformers, so no kinetic or reversibility claim is drawn from them. Polymer chains carrying many grafts, interactions between neighboring grafts, and synthetic routes are out of scope for this study.

Three next steps follow. First, the measured cost breakdown shows where time can be saved without changing results; two such savings are already in the released code, and an MLIP pre-screen could cut the DFT load for larger, more varied libraries. Second, the DMPZ catholytes should be confirmed with the DFT-frequency thermal tier, and λ\_i recomputed in solvent, which the current code already supports. Third, the proposed cells need experimental tests. The code, worked examples, and test suite are released so that the reported numbers can be checked and extended independently.

\section*{Associated Content}

\begin{suppinfo}

The Supporting Information (available free of charge at [TODO: DOI/URL]) contains:
\begin{itemize}
\item S1, the code module behind each pipeline stage;
\item S2, the equations used for the potentials, capacities, reorganization energies, and thermal corrections;
\item S3, every benchmark endpoint at both thermal tiers, with family, wave, and bootstrap statistics and the provenance of the benchmark record;
\item S4, the BMEPZ calibration check;
\item S5, the evidence for the open-shell SCF guess policy, and a positive control (a molecule known to have two competing solutions, used to show the policy catches it);
\item S6, the quality-check counts and the flagged states;
\item S7, the full scored table (potentials, both capacity conventions, Born λ\_o, Pareto-front membership);
\item S8, the measured cost of each DFT state;
\item S9, software versions.
\end{itemize}
Long tables are also provided as CSV files.
\end{suppinfo}

\subsection*{Code availability}

Versioned code and curated metadata are maintained in the public project repository (https://github.com/Ray16/RTFB-materials-research). The release provides the \texttt{redox} command registry and five worked examples under \texttt{examples/} that rebuild stored scientific results without rerunning DFT (checked with \texttt{scripts/validation/check\_refactor\_oracles.py}). It also includes a test suite of 780 passing tests (1 skipped; \texttt{pytest tests -q}, run 2026-11-18) and a hash-checked thermochemistry code that the worked examples depend on. Installation and a locked scientific environment are described by \texttt{setup\_env.sh} and \texttt{docs/INSTALLATION.md}. GPU work is gated through \texttt{gpu\_reserve}; the released tree runs reproduction and tests CPU-only and launches no DFT or Hessian jobs.

\subsection*{Data availability}

The raw quantum-chemical records (geometries, energies, frequencies, and state metadata) for the candidate pool and the benchmark, about 270 GB in total, are kept in a separate data tree that the code reads through a configurable data root. The smaller records behind specific statements in this paper (the BMEPZ check, the SCF-policy control, the benchmark re-scored at the pool's thermal tier, and the throughput measurements) are in the repository under \texttt{manuscript/evidence/} and \texttt{manuscript/tables/}. The raw-record archive does not yet have a public archival DOI; final deposition and author metadata remain to be supplied \textbf{[TODO: author to supply data DOI, repository release tag, author list, contributions, funding, and competing- interest statements]}.

\bibliography{references}
\end{document}
